%! AP Statistics — Unit 1, Lesson 1.5 — Group Activity
\documentclass[10pt]{article}
\usepackage{apstats-article}
\usepackage{apstats-boxes}

%=============================================================================
\begin{document}

\pageheader{Unit 1, Lesson 1.5}{Group Activity: Graphing Quantitative Data}
\namepartnerperiod

\vspace{0.1in}

\begin{tcolorbox}[colback=sky, colframe=navy, boxrule=0.8pt, arc=2mm,
    left=3mm, right=3mm, top=2mm, bottom=2mm]
\small
\textbf{Part 1} (group, 8 min): Read and analyze a dotplot.\\
\textbf{Part 2} (group, 8 min): Complete a back-to-back stemplot and compare two groups.\\
\textbf{Part 3} (group, 5 min): Interpret a histogram from a frequency table.\\
\textbf{Part 4} (individual, 3 min): Reflection.
\end{tcolorbox}

\vspace{0.12in}

% ── PART 1 ───────────────────────────────────────────────────────────────────
{\large\bfseries\color{navy} Part 1 \quad Read a Dotplot}

\vspace{0.1in}

\noindent\small The dotplot below shows the heights (cm) of 15 seedlings measured
two weeks after planting.

\vspace{0.12in}

\begin{center}
\begin{tikzpicture}[scale=1.05]
  \draw[thick] (3.2,0) -- (9.8,0);
  \foreach \x/\lbl in {4/4, 5/5, 6/6, 7/7, 8/8, 9/9} {
    \draw (\x,-0.12) -- (\x,0.12);
    \node[below, font=\small] at (\x,-0.12) {\lbl};
  }
  \node[below, font=\small\itshape] at (6.5,-0.52) {Height (cm)};
  % freq: 4→1, 5→3, 6→4, 7→4, 8→2, 9→1
  \filldraw[navy] (4, 0.38) circle (0.14);
  \foreach \k in {1,...,3} \filldraw[navy] (5, \k*0.38) circle (0.14);
  \foreach \k in {1,...,4} \filldraw[navy] (6, \k*0.38) circle (0.14);
  \foreach \k in {1,...,4} \filldraw[navy] (7, \k*0.38) circle (0.14);
  \foreach \k in {1,...,2} \filldraw[navy] (8, \k*0.38) circle (0.14);
  \filldraw[navy] (9, 0.38) circle (0.14);
\end{tikzpicture}
\end{center}

\vspace{0.12in}

\begin{enumerate}[leftmargin=1.5em, itemsep=8pt]

  \item How many seedlings have a height of exactly 6 cm or 7 cm?
        \blank{2.5cm}
        What percentage of all seedlings does this represent?
        Show your work.\quad\blank{5cm}

  \item What is the range of heights? \blank{3cm}
        Are there any gaps or clusters? Describe what you see.\\[5pt]
        \writeline\\[1pt]\writeline

  \item A student says: ``The distribution looks roughly symmetric.''
        Do you agree? Use evidence from the dotplot.\\[5pt]
        \writeline\\[1pt]\writeline

  \item For this dataset ($n = 15$), is a dotplot a reasonable choice of display?
        What is one limitation it would have if there were 200 seedlings instead?\\[5pt]
        \writeline\\[1pt]\writeline

\end{enumerate}

\vspace{0.12in}

% ── PART 2 ───────────────────────────────────────────────────────────────────
{\large\bfseries\color{navy} Part 2 \quad Complete a Back-to-Back Stemplot}

\vspace{0.1in}

\noindent\small Two AP Statistics classes took the same unit test (100 points). Their scores
are listed below. Complete the back-to-back stemplot and then compare the two distributions.

\vspace{0.1in}

\begin{center}
\small
\begin{tabular}{p{0.42\linewidth} p{0.42\linewidth}}
\textbf{Period 1 ($n = 10$):} & \textbf{Period 2 ($n = 10$):} \\
68, 71, 72, 76, 79, 83, 85, 88, 91, 94 & 65, 69, 72, 75, 77, 78, 82, 84, 88, 91 \\
\end{tabular}
\end{center}

\vspace{0.1in}

\noindent Complete the stemplot below. Leaves for Period 1 go on the \textit{left}
(written right-to-left), leaves for Period 2 go on the \textit{right}.

\vspace{0.1in}

\begin{center}
\renewcommand{\arraystretch}{2.0}
\begin{tabular}{r c l}
\textbf{Period 1} & \textbf{Stem} & \textbf{Period 2} \\
\hline
\blank{3.5cm} & 6 & \blank{3.5cm} \\
\blank{3.5cm} & 7 & \blank{3.5cm} \\
\blank{3.5cm} & 8 & \blank{3.5cm} \\
\blank{3.5cm} & 9 & \blank{3.5cm} \\
\end{tabular}
\hspace{1cm} Key: \textbf{7 $|$ 6} = 76
\end{center}

\vspace{0.1in}

\begin{enumerate}[leftmargin=1.5em, itemsep=8pt]

  \item What is the median score for each class?
        Show how you found each median.\\[5pt]
        Period 1 median: \blank{3cm} \quad Period 2 median: \blank{3cm}

  \item Which class had a higher median? Which class had more spread in scores?
        Support your answer with values from the stemplot.\\[5pt]
        \writeline\\[1pt]\writeline

  \item A student says: ``Period 1 is doing better because it has more students with
        scores in the 80s and 90s.'' Use the back-to-back stemplot to evaluate
        this claim with specific evidence.\\[5pt]
        \writeline\\[1pt]\writeline\\[1pt]\writeline

\end{enumerate}

\vspace{0.12in}

% ── PART 3 ───────────────────────────────────────────────────────────────────
{\large\bfseries\color{navy} Part 3 \quad Interpret a Histogram (from a Frequency Table)}

\vspace{0.1in}

\noindent\small The table below shows the distribution of heights (cm) for a random
sample of 50 adults.

\vspace{0.1in}

\renewcommand{\arraystretch}{1.6}
\begin{tabularx}{\linewidth}{@{} C{0.22\linewidth} C{0.18\linewidth} C{0.22\linewidth} C{0.26\linewidth} @{}}
\rowcolor{sky}
\textbf{Bin} & \textbf{Frequency} & \textbf{Rel.\ Freq.} & \textbf{Cumulative Freq.} \\
\midrule
$[150, 160)$ & 3  & \blank{2.3cm} & \blank{2.3cm} \\
$[160, 165)$ & 8  & \blank{2.3cm} & \blank{2.3cm} \\
$[165, 170)$ & 15 & \blank{2.3cm} & \blank{2.3cm} \\
$[170, 175)$ & 14 & \blank{2.3cm} & \blank{2.3cm} \\
$[175, 180)$ & 7  & \blank{2.3cm} & \blank{2.3cm} \\
$[180, 190)$ & 3  & \blank{2.3cm} & \blank{2.3cm} \\
\midrule
\rowcolor{sky}\textbf{Total} & 50 & \blank{2.3cm} & \blank{2.3cm} \\
\end{tabularx}

\vspace{0.1in}

\begin{enumerate}[leftmargin=1.5em, itemsep=8pt]

  \item Complete the relative frequency and cumulative frequency columns above.

  \item Which bin has the greatest frequency? What percentage of adults fall in
        this bin?\\[5pt]
        \writeline

  \item Notice that the bins are \textit{not} equal width: $[150,160)$ is 10 cm wide
        while $[160,165)$ is only 5 cm wide. Why might this be a problem when
        comparing bar heights in a histogram?\\[5pt]
        \writeline\\[1pt]\writeline

  \item What percentage of adults are at least 170 cm tall?
        Show your calculation.\\[5pt]
        \writeline

\end{enumerate}

\vspace{0.12in}

% ── PART 4 ───────────────────────────────────────────────────────────────────
{\large\bfseries\color{navy} Part 4 \quad Individual Reflection}

\vspace{0.08in}

\begin{tcolorbox}[colback=goldbg, colframe=goldacc, boxrule=0.8pt, arc=2mm,
    left=3mm, right=3mm, top=2mm, bottom=2mm]
\small\textbf{In your own words:} A dotplot, stemplot, and histogram can all display
the same quantitative dataset. What is the most important trade-off to consider
when choosing between them? Give a specific example from today's activity.\\[6pt]
\writeline\\[1pt]\writeline\\[1pt]\writeline
\end{tcolorbox}

\end{document}
